%%%%%%%%%%%%%%%%%%%%%%%%%%%%%%%%%%%%%%%
% Deedy - One Page Two Column Resume
% LaTeX Template
% Version 1.3 (22/9/2018)
%
% Original author:
% Debarghya Das (http://debarghyadas.com)
%
% Original repository:
% https://github.com/deedydas/Deedy-Resume
%
% v1.3 author:
% Zachary Taylor
%
% v1.3 repository:
% https://github.com/ZDTaylor/Deedy-Resume-Reversed
%
% IMPORTANT: THIS TEMPLATE NEEDS TO BE COMPILED WITH XeLaTeX
%
%%%%%%%%%%%%%%%%%%%%%%%%%%%%%%%%%%%%%%

\documentclass[]{deedy-resume-reversed}
\usepackage{fancyhdr}

\pagestyle{fancy}
\fancyhf{}

\begin{document}

%%%%%%%%%%%%%%%%%%%%%%%%%%%%%%%%%%%%%%
%
%     LAST UPDATED DATE
%
%%%%%%%%%%%%%%%%%%%%%%%%%%%%%%%%%%%%%%
% \lastupdated

%%%%%%%%%%%%%%%%%%%%%%%%%%%%%%%%%%%%%%
%
%     TITLE NAME
%
%%%%%%%%%%%%%%%%%%%%%%%%%%%%%%%%%%%%%%
\namesection{Rajesh Kumar}{ 
\begin{center}
    \faEnvelope \hspace{1mm} \href{mailto:rajesh8368568776@gmail.com}{rajesh8368568776@gmail.com} \hspace{8mm} || 
    \faPhone \hspace{1mm} \href{tel:+918368568776}{+91-8368568776} \hspace{8mm} || 
    \faLinkedin \hspace{1mm} \href{https://www.linkedin.com/in/rajesh-kumar-564b8a188/}{LinKedin}
\end{center}
}

%%%%%%%%%%%%%%%%%%%%%%%%%%%%%%%%%%%%%%
%
%     COLUMN ONE
%
%%%%%%%%%%%%%%%%%%%%%%%%%%%%%%%%%%%%%%

\begin{minipage}[t]{0.60\textwidth}

%%%%%%%%%%%%%%%%%%%%%%%%%%%%%%%%%%%%%%
%     EXPERIENCE
%%%%%%%%%%%%%%%%%%%%%%%%%%%%%%%%%%%%%%

\section{Experience}
\runsubsection{Innovaccer}
\descript{|| Data Analyst }
\location{Jan 2021 – Aug 2021 || Noida || }
\vspace{\topsep}
\begin{tightemize}
    \item \textbf{Data Cleaning and Processing}: Used \texttt{pandas}, \texttt{numpy}, and SQL to clean and preprocess large-scale patient data.
\end{tightemize}
\runsubsection{Bare International}
\descript{|| ML Engineer }
\location{join by 2024 || East Mumbai || }
\vspace{\topsep}
\begin{tightemize}
    \item Designed and deployed end-to-end ML pipeline for AI-enabled chatbot system, reducing response latency by 40\% through model optimization.
\end{tightemize}

\runsubsection{KIT College of Engineering Kolhapur, Maharashtra India}
\descript{|| Assistant Professor }
\location{join by July 2025 }
\begin{tightemize}
    \item Assistant Professor At CSE (AIML/DS)
\end{tightemize}
\sectionsep


%%%%%%%%%%%%%%%%%%%%%%%%%%%%%%%%%%%%%%
%     PROJECTS
%%%%%%%%%%%%%%%%%%%%%%%%%%%%%%%%%%%%%%

\section{Projects}
\runsubsection{Land Use Land Cover Mapping from Satellite Images}
Github:// \href{https://github.com/rajesh0305/MSFCN-MASTER}{\bf Link} \\
\descript{| Deep Learning Algorithm for Segmentation}
\location{Jan 2025 – Present | M-tech Final Year Project}  
\textbf{Skills:} Python, TensorFlow, Keras, U-Net, Attention U-Net, Deep Learning, GIS, Remote Sensing, AWS, MLflow, DVC  
\github{GitHub}  
\begin{tightemize}
    \item Developed a \textbf{Deep Learning} model for \textbf{Land Use Land Cover (LULC)} classification using high-resolution \textbf{WorldView-3} satellite images and \textbf{WHDLD} Dataset.
    \item Implemented \textbf{Various Benchmark Model} for precise segmentation of satellite imagery to classify urban, vegetation, water, and barren land regions.
    \item Used \textbf{MLflow} for experiment tracking and \textbf{DVC} for dataset management to ensure reproducibility.
    \item Optimized model training on \textbf{Kaggle TPU} to handle large-scale satellite image datasets efficiently.
    \item Deployed the model on \textbf{AWS (EC2 and S3)} for real-time inference and GIS applications.
\end{tightemize}

\runsubsection{Amazon Fine Food Review System}
Github:// \href{https://github.com/rajesh0305/Amazon_Food_Reviews-main?tab=readme-ov-file}{\bf Link} \\
\descript{| Sentiment Analysis using NLP and Machine Learning}
\location{Jan 2021 – Dec 2021 | B-tech Final Year Project}  
\textbf{Skills:} Python, NLP, Machine Learning, Deep Learning, TensorFlow, Scikit-Learn, LSTM, BERT, FastAPI, AWS    
\begin{tightemize}
    \item Built a \textbf{Sentiment Analysis} system to classify customer reviews from the \textbf{Amazon Fine Food dataset} as \textbf{positive}, \textbf{negative}, or \textbf{neutral}.
    \item Preprocessed text using \textbf{NLTK, spaCy, and TF-IDF} for feature extraction and stopword removal.
    \item Implemented \textbf{Machine Learning} models (Logistic Regression, Random Forest, SVM) and \textbf{Deep Learning} models (LSTM, BERT) to improve sentiment classification accuracy.
    \item Deployed the model using \textbf{FastAPI} and \textbf{AWS (EC2, S3)} for real-time inference.
\end{tightemize}

\end{minipage}
\hfill
\begin{minipage}[t]{0.33\textwidth}

%%%%%%%%%%%%%%%%%%%%%%%%%%%%%%%%%%%%%%
%     EDUCATION
%%%%%%%%%%%%%%%%%%%%%%%%%%%%%%%%%%%%%%

\section{Education}
\subsection{DCRUST}
\descript{Deenabandhu Chotu Ram University Of Science and Technology}
\location{ CGPA: 8.02 / 10.0}
\sectionsep
\subsection{National Institute Of Technology, Suratkal}
\descript{Master's Of Technology In Signal Processing And Machine Learning}
\location{ CGPA: 8.45 / 10.0}
\sectionsep

%%%%%%%%%%%%%%%%%%%%%%%%%%%%%%%%%%%%%%
%     SKILLS
%%%%%%%%%%%%%%%%%%%%%%%%%%%%%%%%%%%%%%

\section{Skills}  
\textbf{Programming:} Python, C, Java \\ 
\textbf{Framework:} Streamlit,Flask,FastAPI\\
\textbf{Cloud and Deployment:} AWS (EC2, S3), Docker, MLflow, DVC.\\
\textbf{ML/DL:} TensorFlow, PyTorch, Scikit-Learn, NLP, Computer Vision, Transformers, XGBoost ,Supervised and Unsupervised Machine Learning Algorithms, Hyperparameter Tuning, Feature Selection, EDA.\\  
\textbf{Data Analysis \& Visualization:} pandas, NumPy, Matplotlib, Seaborn.\\
\textbf{Databases} MySQL,SQL \\  
\textbf{NLP \& LLM:}LSTM, Transformer, Bert, Langchain, Agents, RAG.\\
\textbf{Statistics:}Hypothesis testing, Descriptive analysis.
\sectionsep

%%%%%%%%%%%%%%%%%%%%%%%%%%%%%%%%%%%%%%
%     COURSEWORK
%%%%%%%%%%%%%%%%%%%%%%%%%%%%%%%%%%%%%%

\section{Coursework}
\subsection{}
Probability and statistics, Linear Algebra, Machine Learning and Deep Learning, Database Management System, Data Structures and Algorithms, Object-Oriented-programing, Optimization, Operating system.
\sectionsep

%%%%%%%%%%%%%%%%%%%%%%%%%%%%%%%%%%%%%%
%     LINKS
%%%%%%%%%%%%%%%%%%%%%%%%%%%%%%%%%%%%%%

\section{Links}
Github:// \href{https://github.com/rajesh0305/}{\bf Rajesh Kumar} \\
LinkedIn://  \href{https://www.linkedin.com/in/rajesh-kumar-564b8a188/}{\bf Rajesh Kumar}\\
CODE:// \href{https://www.geeksforgeeks.org/user/rajeshrajput/}{\bf GFG}\\
CODE:// \href{https://leetcode.com/u/datafunda/}{\bf LeetCode}
\sectionsep

\end{minipage}

% ====================== SECOND PAGE (EXACT IMAGE FORMAT) ======================
\newpage
\pagestyle{fancy}
\fancyhf{}
\renewcommand{\headrulewidth}{0pt}
\cfoot{\thepage}

% Header (same as first page)
\namesection{Rajesh Kumar}{ 
\begin{center}
    \faEnvelope \hspace{1mm} \href{mailto:rajesh8368568776@gmail.com}{rajesh8368568776@gmail.com} \hspace{8mm} || 
    \faPhone \hspace{1mm} \href{tel:+918368568776}{+91-8368568776} \hspace{8mm} || 
    \faLinkedin \hspace{1mm} \href{https://www.linkedin.com/in/rajesh-kumar-564b8a188/}{LinkedIn}
\end{center}
}

% Use the resume's standard section style
\section{Research Coordinator at KIT College of Engineering, Kolhapur}

% Role and date (using same style as Experience entries)
\runsubsection{AIML/DS Department}
\descript{|| Research Coordinator (Proposed Role)}
\location{Joining: July 2025}

\vspace{\topsep}
\begin{tightemize}
    \item \textbf{Overview:} Foster a vibrant research ecosystem bridging academia and industry. Elevate departmental research output, secure funding, and mentor students/faculty in AI/ML/DL.
\end{tightemize}

% --- KEY ---
\subsection*{KEY}//
\begin{tightemize}
    \item \textbf{Strategic Research Planning:} Develop a 3-year research roadmap aligned with national/global aspirations.
    \item \textbf{Faculty \& Student Mentorship:} Lead weekly seminars/workshops; guide M.Tech and Ph.D. scholars.
    \item \textbf{Collaboration \& Partnerships:} Forge links with industry (Tech Mahindra, TCS, Accenture) and academia (IIITs, NITs).
    \item \textbf{Funding \& Grants:} Write proposals for DST, SERB, AICTE, CSR to secure infrastructure/project funding.
    \item \textbf{Research Output \& Visibility:} Drive high-impact publications (SCI/Scopus); manage repository and online presence.
    \item \textbf{Infrastructure Development:} Set up state-of-the-art labs for Deep Learning, NLP, Computer Vision with cloud resources.
\end{tightemize}

% --- FUTURE ---
\subsection*{FUTURE}//
\begin{tightemize}
    \item \textbf{Center of Excellence (CoE) in AI/ML:} Establish KIT as a premier AI research hub with groups in Healthcare, Agriculture, and NLP.
    \item \textbf{Industry-Academia Integration (Sandwich Model):} Launch a Research Internship Program (6–12 months on industry problems) and quarterly Industry-Academia Conclaves.
    \item \textbf{International Collaborations:} MoUs with foreign universities; invite global experts as Visiting Professors.
    \item \textbf{Startup \& Innovation Incubation:} Create an AI/ML Incubation Cell; organise annual hackathons/ideathons for local problems.
\end{tightemize}

% --- PROPOSED (First-Year Timeline) ---
\subsection*{PROPOSED}//
\begin{tightemize}
    \item \textbf{Months 1–3 (Jul–Sep 2025):} Onboarding, finalise roadmap, assess labs, meet faculty/students.
    \item \textbf{Months 4–6 (Oct–Dec 2025):} Start research seminars, draft funding proposals, reach out to industry.
    \item \textbf{Months 7–9 (Jan–Mar 2026):} First Industry-Academia Conclave, begin MoU process, call for student projects.
    \item \textbf{Months 10–12 (Apr–Jun 2026):} Review progress, measure KPIs, plan Year 2, host a mini-symposium.
\end{tightemize}

\vspace{0.3cm}
\textit{Note: This plan is dynamic and will be refined through consultations with faculty, management, and industry partners.}

% ====================== END OF SECOND PAGE ======================

\end{document}